\section{Invariant-subspace decomposition and irreducible blocks}

The block decomposition rests on one principle: a positive semidefinite fixed
point of $\E_A$ that is not full rank has a proper invariant support, which
block-diagonalizes the Kraus operators without changing the MPV family. Iterating
the resulting splitting decomposes any tensor into irreducible blocks.

\begin{theorem}[Unitary conjugation preserves the MPV family]\label{thm:samempv_unitary}
    \lean{MPSTensor.sameMPV_conj_unitary}
    \leanok
    \uses{def:same_mpv}
    If $U$ is a unitary matrix, then $A$ and $U^\dagger A U$
    generate the same MPV family.
\end{theorem}

\begin{proof}\leanok
    By cyclicity of the trace,
    $\tr((U^\dagger A^{i_1} U) \cdots (U^\dagger A^{i_N} U))
        = \tr(A^{i_1} \cdots A^{i_N})$.
\end{proof}

\begin{theorem}[Invariant projection gives a two-block decomposition]\label{thm:lowerzero_two_block}
    \lean{MPSTensor.exists_twoBlock_decomp_of_lowerZero}
    \leanok
    \uses{def:same_mpv2}
    Suppose $P$ is an orthogonal projection satisfying, for every physical
    index~$i$,
    \[
        (\Id - P) A^i P = 0.
    \]
    Then there exist $n,m$ with $n + m = D$ and MPS tensors $A_1, A_2$ of bond
    dimensions $n$ and $m$ such that $A$ and the two-block tensor
    $A_1 \oplus A_2$ generate the same MPV family in the sense of
    Definition~\ref{def:same_mpv2}.
\end{theorem}

\begin{proof}
    \leanok
    \uses{thm:samempv_unitary}
    Diagonalize $P$ by a unitary $U$.
    Since $P^2 = P$, the diagonal form of $P$ has only $0$- and $1$-eigenvalues,
    so the basis splits as $n + m = D$.
    Conjugating $A$ by $U$ preserves the MPV family by
    Theorem~\ref{thm:samempv_unitary}.
    In this basis the relation $(\Id - P) A^i P = 0$ makes each letter upper
    block triangular.
    Write
    \[
        B^i =
        \begin{pmatrix}
            A_1^i & *     \\
            0     & A_2^i
        \end{pmatrix}
    \]
    for the conjugated letters in this split basis. Since
    $B^{i_1}\cdots B^{i_N}$ is again upper block triangular, its trace satisfies
    \[
        \tr(B^{i_1}\cdots B^{i_N})
        =
        \tr(A_1^{i_1}\cdots A_1^{i_N})
        +
        \tr(A_2^{i_1}\cdots A_2^{i_N}).
    \]
    Hence the block-diagonal tensor $A_1 \oplus A_2$ has the same MPV
    coefficients as the conjugated tensor.
\end{proof}

\begin{theorem}[Nontrivial invariant projection gives strict two-block reduction]
    \label{thm:lowerzero_two_block_strict}
    \lean{MPSTensor.exists_twoBlock_decomp_of_lowerZero_strict}
    \leanok
    \uses{def:same_mpv2}
    Under the hypotheses of Theorem~\ref{thm:lowerzero_two_block},
    assume in addition that $P \neq 0$ and $P \neq \Id$.
    Then there exist $n,m$ with $n + m = D$, $n < D$, and $m < D$, together
    with block tensors $A_1, A_2$, such that
    $A$ and $A_1 \oplus A_2$ generate the same MPV family in the sense of
    Definition~\ref{def:same_mpv2}.
\end{theorem}

\begin{proof}
    \leanok
    \uses{thm:lowerzero_two_block}
    Apply Theorem~\ref{thm:lowerzero_two_block}.
    In the diagonal form of $P$, the $1$-eigenspace is nonempty because
    $P \neq 0$, and the $0$-eigenspace is nonempty because $P \neq \Id$.
    Hence both block sizes are positive.
    Since $n + m = D$, this implies $n < D$ and $m < D$.
\end{proof}

\begin{definition}[Support projection]\label{def:support_proj}
    \lean{MPSTensor.supportProj}
    \leanok
    For a positive semidefinite matrix $\rho \ge 0$, the \emph{support
        projection} $P$ is the orthogonal projection onto the range of $\rho$. Via
    the spectral decomposition
    $\rho = U \diag(\lambda_1,\ldots,\lambda_D) U^\dagger$,
    \[
        P = U \diag(\mathbf{1}_{\lambda_1>0},\ldots,\mathbf{1}_{\lambda_D>0}) U^\dagger,
    \]
    where $\mathbf{1}_{\lambda_j>0}$ is $1$ if $\lambda_j > 0$ and $0$ otherwise.
\end{definition}

\begin{definition}[Invariant projection predicate]\label{def:has_inv_proj}
    \lean{MPSTensor.HasInvariantProj}
    \leanok
    \uses{def:mps_tensor}
    An MPS tensor $A$ \emph{has an invariant projection} if there exists an
    orthogonal projection $P$ with $P \neq 0$, $P \neq \Id$, and
    $(\Id - P) A^i P = 0$ for all~$i$. A tensor is \emph{irreducible} if it has
    no nontrivial invariant projection (Definition~\ref{def:irreducible_tensor}).
\end{definition}

\begin{theorem}[PSD fixed point gives invariant projection]\label{thm:fp_inv_proj}
    \lean{MPSTensor.lowerZero_of_posSemidef_fixedPoint}
    \leanok
    \uses{def:transfer_map}
    % Source: the support-projection step in the proof of Theorem 4
    % (PerezGarcia2007Matrix): the support of a singular positive fixed point is
    % invariant, A_i P_R = P_R A_i P_R, by positivity.
    Let $\rho \ge 0$ be a fixed point of $\E_A$, and let $P$ be its support
    projection. Then $(\Id - P) A^i P = 0$ for all~$i$: the range of $P$ is
    invariant under every $A^i$.
\end{theorem}

\begin{proof}
    \leanok
    From $\E_A(\rho) = \rho$ and positivity of each summand
    $A^i \rho (A^i)^\dagger$, one gets
    $(\Id - P) A^i \rho (A^i)^\dagger (\Id - P) = 0$ for each~$i$. Since $\rho$
    is supported on the range of $P$, this forces $(\Id - P) A^i P = 0$.

    This is the finite-dimensional case of~\cite[Proposition~6.10]{Wolf2012Quantum}:
    the support projection of any stationary state is sub-harmonic. The
    equivalence between irreducibility and absence of nontrivial invariant
    projections is~\cite[Theorem~6.2]{Wolf2012Quantum}; see
    also~\cite[Section~IV.A.1]{Cirac2021Matrix}.
\end{proof}

\begin{theorem}[Two-block decomposition]\label{thm:two_block_decomp}
    \lean{MPSTensor.exists_twoBlock_decomp_of_posSemidef_fixedPoint_strict}
    \leanok
    \uses{def:transfer_map, def:same_mpv2}
    % Source: the invariant-support projection and finite-ring trace split in
    % the proof of Theorem 4 (PerezGarcia2007Matrix).
    If $\E_A$ has a nonzero PSD fixed point $\rho$ that is not positive definite,
    then there exist MPS tensors $A_1, A_2$ of smaller bond dimensions with
    $\mc{V}(A) = \mc{V}(A_1 \oplus A_2)$.
\end{theorem}

\begin{proof}
    \leanok
    \uses{thm:fp_inv_proj, thm:samempv_unitary}
    The support projection $P$ of $\rho$ is neither $0$ nor $\Id$ (since
    $\rho \neq 0$ and $\rho$ is not positive definite).
    Theorem~\ref{thm:fp_inv_proj} gives $(\Id - P) A^i P = 0$, so a unitary
    change of basis block-diagonalizing $P$ puts $A^i$ in upper-triangular block
    form with diagonal blocks $B_{11}^i$, $B_{22}^i$.
    Theorem~\ref{thm:samempv_unitary} preserves the MPV under the conjugation,
    and the trace of an upper-triangular product equals the sum of the
    diagonal-block traces, so deleting the off-diagonal blocks preserves all MPV
    coefficients.
\end{proof}

\subsection{Iterated reduction to irreducible blocks}

Iterating the two-block splitting on bond dimension decomposes any tensor into
irreducible blocks; the remaining lemmas record the spectral facts about
irreducible unital blocks and the algebra-spanning input used later.

\begin{definition}[Irreducible tensor]\label{def:irreducible_tensor}
    \lean{MPSTensor.IsIrreducibleTensor}
    \leanok
    \uses{def:has_inv_proj}
    An MPS tensor~$A$ is \emph{irreducible} if it has no nontrivial invariant
    projection.
\end{definition}

\begin{theorem}[Irreducible block decomposition]\label{thm:irred_decomp}
    \lean{MPSTensor.exists_irreducible_blockDecomp}
    \leanok
    \uses{def:irreducible_tensor}
    Every MPS tensor $A$ admits a block-diagonal tensor $\bigoplus_{k=1}^r A_k$
    with each $A_k$ irreducible,
    $\mc{V}(A) = \mc{V}(\bigoplus_k A_k)$, and
    $\sum_{k=1}^rD_k=D$.
\end{theorem}

\begin{proof}
    \leanok
    \uses{thm:two_block_decomp}
    By strong induction on $D$. If $A$ is irreducible, take $r = 1$. Otherwise a
    nontrivial invariant projection $P$ gives the upper-triangular splitting of
    Theorem~\ref{thm:two_block_decomp}; deleting the off-diagonal part preserves
    the MPV family and yields two blocks of strictly smaller bond dimension.
    Apply the induction hypothesis to each block.

    Both~\cite[Theorem~4]{PerezGarcia2007Matrix}
    and~\cite[Section~IV.A.1]{Cirac2021Matrix} obtain the block decomposition by
    the same invariant-projection splitting.
\end{proof}

\section{Nonzero blocks and the trace-preserving gauge}%
\label{sec:nonzero_blocks_tp_gauge}

In an irreducible block decomposition some blocks may have all-zero Kraus
operators. Such a block contributes nothing to a matrix product vector at any
positive length and is therefore discarded. The structural direct-sum identity
implies that the retained block dimensions satisfy $\sum_k D_k\le D$ without
using an empty-word coefficient. Each retained block is then placed in the
trace-preserving gauge.

\begin{theorem}[Nonzero irreducible block decomposition]%
    \label{thm:nonzero_irreducible_block_decomposition}
    \lean{MPSTensor.exists_irreducible_blockDecomp_nonzeroBlocks}
    \leanok
    \uses{def:irreducible_tensor, def:same_mpv2_pos}
    Every MPS tensor $A$ admits a finite family of nonzero irreducible blocks
    $(B_k)_{k=1}^r$, each with at least one nonzero Kraus operator and positive
    bond dimension. For every $N \ge 1$ and every $\sigma$,
    \[
        V^{(N)}(A)_\sigma
        = V^{(N)}\!(\textstyle\bigoplus_{k=1}^r B_k)_\sigma.
    \]
    Moreover, $\sum_{k=1}^r D_k\le D$.
\end{theorem}

\begin{proof}\leanok
    \uses{thm:irred_decomp}
    Apply the irreducible block decomposition (Theorem~\ref{thm:irred_decomp})
    and retain the blocks having a nonzero Kraus operator. For every all-zero
    irreducible block $C_j$ and every $N \ge 1$, $V^{(N)}(C_j)_\sigma = 0$, so
    \[
        V^{(N)}(A)_\sigma
        = V^{(N)}\!(\textstyle\bigoplus_{k=1}^r B_k)_\sigma.
    \]
    The dimension inequality follows by restricting the structural identity
    $\sum_j D_j=D$ to the retained subset.
\end{proof}

\begin{theorem}[Blockwise TP-gauge normalization of irreducible blocks]%
    \label{thm:blockwise_tp_gauge}
    \lean{MPSTensor.exists_tp_gauge_blockwise}
    \leanok
    \uses{def:irreducible_tensor}
    Suppose $A$ generates the same MPV family as $\bigoplus_{k=1}^r B_k$, where
    each $B_k$ is irreducible with at least one nonzero Kraus operator. Then there
    are nonzero weights $(\mu_k)_{k=1}^r$ and blocks $(C_k)_{k=1}^r$ such that, for
    every $N$ and every $\sigma$,
    \[
        V^{(N)}(A)_\sigma
        =
        V^{(N)}\!(\textstyle\bigoplus_{k=1}^r \mu_k C_k)_\sigma
        =
        \sum_{k=1}^r \mu_k^{\,N} V^{(N)}(C_k)_\sigma,
    \]
    each $C_k$ irreducible and left-canonical,
    $\sum_i (C_k^i)^\dagger C_k^i = \Id$, with positive bond dimension. This is
    the blockwise form used in the arbitrary-tensor trace-preserving gauge
    below.
\end{theorem}

\begin{proof}\leanok
    The Perron--Frobenius TP-gauge construction is applied to each irreducible
    block. Gauge equivalence preserves irreducibility and the MPV family, and the
    positive scaling factors are absorbed into the weights $\mu_k$.
\end{proof}

\subsection{Trace-preserving gauge for arbitrary tensors}%
\label{sec:tp_gauge_arbitrary}

\begin{theorem}[TP-gauge reduction for arbitrary tensors]%
    \label{thm:tp_gauge_arbitrary}
    \lean{MPSTensor.exists_tp_gauge_from_arbitrary}
    \leanok
    \uses{def:irreducible_tensor, def:same_mpv2_pos}
    For any MPS tensor $A$, there exist nontrivial blocks $(B_k)_{k=1}^r$ with
    nonzero weights $(\mu_k)_{k=1}^r$,
    each block irreducible, left-canonical
    ($\sum_i (B_k^i)^\dagger B_k^i = \Id$), and of positive bond dimension, such
    that for every $N \ge 1$ and every $\sigma$,
    \[
        V^{(N)}(A)_\sigma
        = V^{(N)}\!(\textstyle\bigoplus_{k=1}^r \mu_k B_k)_\sigma,
    \]
    and $\sum_{k=1}^r D_k\le D$.
\end{theorem}

\begin{proof}\leanok
    \uses{thm:nonzero_irreducible_block_decomposition, thm:tp_data_irreducible}
    Theorem~\ref{thm:nonzero_irreducible_block_decomposition} supplies the nonzero
    irreducible blocks $(C_k)_{k=1}^r$. For every $N \ge 1$ and every $\sigma$,
    \[
        V^{(N)}(A)_\sigma
        = V^{(N)}\!(\textstyle\bigoplus_{k=1}^r C_k)_\sigma.
    \]
    The blockwise TP-gauge theorem then replaces each $C_k$ by a
    left-canonical block $B_k$ with weight $\mu_k$. Again for every $N \ge 1$
    and every $\sigma$,
    \[
        V^{(N)}\!(\textstyle\bigoplus_{k=1}^r C_k)_\sigma
        =
        V^{(N)}\!(\textstyle\bigoplus_{k=1}^r \mu_k B_k)_\sigma.
    \]
    The gauge preserves every block dimension, so the structural bound
    $\sum_kD_k\le D$ is unchanged.
\end{proof}

\section{Blocking, period removal, and primitive blocks}%
\label{sec:blocking_infrastructure}

Blocking by an integer~$p$ carries MPV equalities, weights, transfer maps, and
identifications of the blocked physical alphabet. The first consequence is the
existence of a common blocking period under which every transfer map in a finite
family becomes primitive; the period-removal and primitive-block reductions then
proceed from that common period.

\begin{theorem}[Common blocking period for a block family]%
    \label{thm:common_blocking_period}
    \lean{MPSTensor.exists_common_blocking_all_primitive}
    \leanok
    \uses{def:blocked_tensor}
    Let $(B_k)_{k=1}^r$ be a finite family of positive-bond-dimension MPS
    tensors, each admitting a blocking period $p_k$ making its transfer map
    primitive. Then $p = \lcm(p_1, \ldots, p_r)$ is a common
    period: $\E_{B_k^{[p]}}$ is primitive for every~$k$.
\end{theorem}

\begin{proof}\leanok
    Each $p_k$ divides $p$, so primitivity of $\E_{B_k^{[p_k]}}$ carries over to
    $\E_{B_k^{[p]}}$ via the divisibility identity for transfer-map powers.
\end{proof}

The following stability statements are used when a primitive irreducible block
is blocked again for a common-refinement or injectivity length. This later
blocking is distinct from the period-removal length that produced the sector
block.

\begin{theorem}[Blocking preserves irreducibility for primitive irreducible blocks]%
    \label{thm:blocked_irreducible_tp_primitive}
    \lean{MPSTensor.isIrreducibleTensor_blockTensor_of_tp_primitive_irr}
    \leanok
    \uses{def:blocked_tensor, def:irreducible_tensor}
    Let $A$ be left-canonical, irreducible, primitive, of positive bond
    dimension. For every $P>0$, the blocked tensor $A^{[P]}$ is irreducible:
    there is no orthogonal projection $Q$ with $Q \ne 0$, $Q \ne \Id$ such that,
    for every $P$-blocked index $\alpha$,
    \[
        (\Id-Q)(A^{[P]})^\alpha Q=0.
    \]
\end{theorem}

\begin{proof}\leanok
    The blocked transfer map is the corresponding positive power of the original
    one:
    \[
        \E_{A^{[P]}}(X)=\E_A^P(X).
    \]
    Hence a positive-definite fixed point of $\E_A$ remains fixed by
    $\E_{A^{[P]}}$. Primitivity gives uniqueness of positive semidefinite fixed
    points for the blocked map, hence its transfer map is irreducible, which is
    the irreducibility criterion for $A^{[P]}$.
\end{proof}

\begin{theorem}[Extra blocking of primitive irreducible sector blocks]%
    \label{thm:tp_primitive_irreducible_extra_blocking}
    \lean{MPSTensor.tp_primitive_irreducible_extra_blocking}
    \leanok
    \uses{def:blocked_tensor, thm:blocked_irreducible_tp_primitive}
    Let $A$ be trace-preserving, primitive, and tensor-irreducible. For every
    $k>0$,
    \[
        \sum_\alpha ((A^{[k]})^\alpha)^\dagger
        (A^{[k]})^\alpha = \Id,
        \qquad
        \sigma_\partial(\E_{A^{[k]}})=\{1\},
    \]
    and $A^{[k]}$ is tensor-irreducible. This $k$ is a later
    common-refinement or injectivity length, distinct from the period-removal
    length that produced the sector block.
\end{theorem}

\begin{proof}\leanok
    \uses{thm:blocked_irreducible_tp_primitive}
    Trace preservation follows from blocking a trace-preserving tensor,
    primitivity from a positive power of a primitive transfer map, and
    irreducibility from Theorem~\ref{thm:blocked_irreducible_tp_primitive}.
\end{proof}

\subsection{Period removal by cyclic sectors}%
\label{sec:cyclic_sectors}

The step needed below is the period-removing blocking described
in~\cite[Section~2.3]{Cirac2016MPDO_arXiv}: after blocking each irreducible TP
block by the least common multiple of its peripheral periods, the result has no
nontrivial $p$-periodic vectors. The cyclic-sector
decomposition comes from the peripheral spectrum of the adjoint transfer
map~\cite[Theorem~6.6]{Wolf2012Quantum}.

\begin{theorem}[Period removal by blocking]%
    \label{thm:cyclic_sector_decomp_after_blocking}
    \lean{MPSTensor.exists_cyclic_sector_decomp_after_blocking_of_TP_of_isIrreducibleTensor}
    \leanok
    \uses{def:blocked_tensor, thm:peripheral_cyclic_structure}
    Let $A$ be a left-canonical irreducible tensor. Then there is a
    period-removing length $m\ge 1$ such that $A^{[m]}$ admits:
    \begin{enumerate}
        \item left-canonical sector tensors $(C_k)_{k=1}^m$ of nonzero bond
              dimension,
        \item a unit-weight direct-sum MPV decomposition
              $A^{[m]} \sim \bigoplus_{k=1}^m C_k$,
        \item orthogonal projections $(P_k)_{k=1}^m$ summing to $\Id$ with
              $\E_{A^\dagger}(P_{k+1}) = P_k$,
        \item commutation relations
              $P_k A^{[m]}_i = A^{[m]}_i P_k$ for every blocked letter $i$,
        \item the sector trace formula
              $f_{C_k}(\sigma)=
                  \tr(P_k A^{[m]}_{\sigma_1}\cdots A^{[m]}_{\sigma_N})$,
        \item linear isomorphisms onto the compression corners
              $\varphi_k : \MN{\dim C_k} \xrightarrow{\sim} P_k \MN{D} P_k$,
        \item transfer-intertwining identities
              $\varphi_k(\E_{C_k^\dagger}(X)) =
                  \E_{(P_kA^{[m]})^\dagger}(\varphi_k(X))$,
        \item multiplicativity and adjoint compatibility
              $\varphi_k(XY)=\varphi_k(X)\varphi_k(Y)$ and
              $\varphi_k(X^\dagger)=\varphi_k(X)^\dagger$.
    \end{enumerate}
    This is the cyclic-sector part of the period-removal step
    of~\cite[Section~2.3]{Cirac2016MPDO_arXiv}; the projections $P_k$ come
    from the peripheral spectrum of the adjoint transfer map
    (\cite[Theorem~6.6]{Wolf2012Quantum}).
\end{theorem}

\begin{proof}\leanok
    The cyclic decomposition of the adjoint transfer map gives projections
    $(P_k)$ with $\E_{A^\dagger}(P_{k+1}) = P_k$. Iterating the cyclic relation
    $m$ times gives $(\E_{A^\dagger})^m(P_k) = P_k$, and the blocked-transfer
    identity $\E_{(A^{[m]})^\dagger} = (\E_{A^\dagger})^m$ identifies these with
    the adjoint transfer map of $A^{[m]}$. Apply the block decomposition from
    adjoint-fixed projections.
\end{proof}

\begin{theorem}[Unconditional primitive blocked sectors]%
    \label{thm:blocked_sector_unconditional}
    \lean{MPSTensor.primitive_and_irreducible_sectorBlocks_of_cyclic_decomp_after_blocking}
    \leanok
    \uses{thm:cyclic_sector_decomp_after_blocking,
        thm:sector_irred_unconditional,
        thm:primitive_adjoint_equiv}
    Let $A$ be a left-canonical irreducible tensor whose adjoint transfer map has
    peripheral spectrum $\{\gamma^j : 0 \le j < m\}$ for a primitive $m$th root of
    unity $\gamma$, with blocked cyclic-sector decomposition $(C_u)_{u=1}^m$,
    $(P_u)_{u=1}^m$, and
    linear isomorphisms
    $\varphi_u : \MN{\dim C_u} \xrightarrow{\sim} P_u \MN{D} P_u$. Assume
    explicitly that the $P_u$ are orthogonal projections summing to $\Id$ and
    satisfying $\E_{A^\dagger}(P_{u+1})=P_u$, that each sector has nonzero bond
    dimension, and that the compression maps satisfy
    \[
        \varphi_u(\E_{C_u^\dagger}(X))
        =
        \E_{(P_uA^{[m]})^\dagger}(\varphi_u(X)),
        \qquad
        \varphi_u(XY)=\varphi_u(X)\varphi_u(Y),
        \qquad
        \varphi_u(X^\dagger)=\varphi_u(X)^\dagger.
    \]
    Here $P_uA^{[m]}$ denotes the blocked tensor with letters
    $P_u(A^{[m]})^\alpha$.
    Then every
    sector tensor satisfies $\sigma_\partial(\E_{C_u})=\{1\}$ and is
    tensor-irreducible.
\end{theorem}

\begin{proof}\leanok
    \uses{thm:sector_irred_unconditional,
        thm:primitive_adjoint_equiv}
    Theorem~\ref{thm:sector_irred_unconditional} gives irreducibility of
    $(\E_A^\dagger)^m$ on each corner $P_u$. The compression equivalences
    $\varphi_u$ intertwine the adjoint transfer map of $C_u$ with the
    corresponding corner restriction of $(\E_A^\dagger)^m$. The peripheral
    spectrum of the blocked adjoint map is
    \[
        \sigma_\partial((\E_A^\dagger)^m)
        =
        \{(\gamma^j)^m : 0 \le j < m\}
        =
        \{1\},
    \]
    because $\gamma$ is a primitive $m$th root of unity. Hence each nonzero
    irreducible corner restriction is primitive. Primitivity and irreducibility
    are preserved across $\varphi_u$; then use
    Theorem~\ref{thm:primitive_adjoint_equiv} to pass from the adjoint blocked
    map to the primal transfer map of $C_u$.
\end{proof}

\begin{theorem}[Period removal with primitive irreducible sectors]%
    \label{thm:cyclic_sector_decomp_irr_tp_prim_irr}
    \lean{MPSTensor.exists_primitive_irreducible_cyclic_sector_decomp_of_TP_of_isIrreducibleTensor}
    \leanok
    \uses{thm:cyclic_sector_decomp_after_blocking, thm:blocked_sector_unconditional}
    Let $A$ be a left-canonical irreducible tensor. Then there is a period
    $m \ge 1$ and a unit-weight decomposition of $A^{[m]}$ into sector blocks
    $(C_k)_{k=1}^m$, each of which is left-canonical, has primitive transfer map,
    is tensor-irreducible, and has nonzero bond dimension. Equivalently, the
    cyclic-sector decomposition of $A^{[m]}$ already satisfies the
    trace-preservation, primitivity, and irreducibility hypotheses needed for the
    common-blocking step
    (Lemma~\ref{thm:exists_common_blocked_cyclic_sector_family}).
\end{theorem}

\begin{proof}\leanok
    \uses{thm:peripheral_cyclic_structure,
        thm:cyclic_sector_decomp_after_blocking, thm:blocked_sector_unconditional}
    Left-canonicality makes $K_i := (A^i)^\dagger$ unital, irreducibility passes
    to the adjoint map, and the quantum Perron--Frobenius theorem provides a
    positive-definite fixed point; the cyclic peripheral-spectrum theorem and
    Theorem~\ref{thm:cyclic_sector_decomp_after_blocking} then produce the
    cyclic-sector decomposition. Keep the projections and compression
    equivalences, and apply
    Theorem~\ref{thm:blocked_sector_unconditional} to the sector blocks.
\end{proof}

\begin{theorem}[Cyclic sectors after removing zero blocks]%
    \label{thm:ft_after_blocking_per_block_cyclic_nonzero_blocks}
    \lean{MPSTensor.afterBlocking_perBlockCyclicData_of_sameMPV₂Pos}
    \leanok
    \uses{thm:tp_gauge_arbitrary, def:same_mpv2_pos,
        def:primitive_irreducible_cyclic_sectors,
        thm:cyclic_sector_decomp_irr_tp_prim_irr}
    If $A$ and $B$ generate the same MPV family at every positive length, then after removing the
    all-zero blocks and applying the trace-preserving gauge there are weighted nonzero-block tensors
    $\bigoplus_j \mu^A_j A_j$ and $\bigoplus_k \mu^B_k B_k$ with, for every
    $N\ge 1$ and $\sigma$,
    \[
        V^{(N)}(A)_\sigma
        =
        V^{(N)}\!(\textstyle\bigoplus_j \mu^A_j A_j)_\sigma,
        \qquad
        V^{(N)}(B)_\sigma
        =
        V^{(N)}\!(\textstyle\bigoplus_k \mu^B_k B_k)_\sigma,
    \]
    and the two nonzero parts agree:
    \[
        V^{(N)}\!(\textstyle\bigoplus_j \mu^A_j A_j)_\sigma
        =
        V^{(N)}\!(\textstyle\bigoplus_k \mu^B_k B_k)_\sigma.
    \]
    Each nonzero block has primitive irreducible cyclic sectors: for each $j$,
    there is $m_j\ge 1$ with a unit-weight sector decomposition
    $A_j^{[m_j]} \sim \bigoplus_{u=1}^{m_j} C^A_{j,u}$, and similarly for $B_k$.
\end{theorem}

\begin{proof}\leanok
    \uses{thm:tp_gauge_arbitrary,
        thm:cyclic_sector_decomp_irr_tp_prim_irr}
    Apply Theorem~\ref{thm:tp_gauge_arbitrary} separately to $A$ and $B$. Its
    positive-length conclusion gives, for $N\ge 1$,
    \[
        V^{(N)}(A)_\sigma
        =
        V^{(N)}\!(\textstyle\bigoplus_j \mu^A_j A_j)_\sigma,
        \qquad
        V^{(N)}(B)_\sigma
        =
        V^{(N)}\!(\textstyle\bigoplus_k \mu^B_k B_k)_\sigma.
    \]
    Chaining these through $V^+(A)=V^+(B)$ identifies the two nonzero parts. Finally
    apply Theorem~\ref{thm:cyclic_sector_decomp_irr_tp_prim_irr} to
    every trace-preserving irreducible nonzero-weight block.
\end{proof}

\begin{theorem}[Common blocking length for cyclic sectors]%
    \label{thm:ft_after_blocking_common_length_common_sector_theorem}
    \lean{MPSTensor.afterBlocking_commonLengthCommonSectorData_of_sameMPV₂Pos}
    \leanok
    \uses{thm:ft_after_blocking_per_block_cyclic_nonzero_blocks, def:same_mpv2_pos,
        thm:exists_common_blocked_cyclic_sector_family_common_multiple,
        thm:same_mpv2_pos_block_tensor_to_tensor_from_blocks,
        thm:same_mpv2_pos_to_tensor_from_blocks_block_power,
        thm:common_blocked_cyclic_sector_derived_properties}
    If two tensors generate the same MPV family at every positive length, then the cyclic-sector
    decompositions on the two sides can be expressed at one common positive
    blocking length $p$. For every $N\ge 1$ and every blocked $\sigma$, each
    blocked tensor equals its powered nonzero part,
    \[
        V^{(N)}\!(A^{[p]})_\sigma
        =
        V^{(N)}\!(\textstyle\bigoplus_j (\mu^A_j)^p (A_j)^{[p]})_\sigma,
        \qquad
        V^{(N)}\!(B^{[p]})_\sigma
        =
        V^{(N)}\!(\textstyle\bigoplus_k (\mu^B_k)^p (B_k)^{[p]})_\sigma,
    \]
    and the two powered nonzero parts agree,
    \[
        V^{(N)}\!\bigl(\textstyle\bigoplus_j
        (\mu^A_j)^p (A_j)^{[p]}\bigr)_\sigma
        =
        V^{(N)}\!\bigl(\textstyle\bigoplus_k
        (\mu^B_k)^p (B_k)^{[p]}\bigr)_\sigma.
    \]
\end{theorem}

\begin{proof}\leanok
    \uses{thm:ft_after_blocking_per_block_cyclic_nonzero_blocks,
        thm:exists_common_blocked_cyclic_sector_family_common_multiple,
        thm:same_mpv2_pos_block_tensor_to_tensor_from_blocks,
        thm:same_mpv2_pos_to_tensor_from_blocks_block_power,
        thm:common_blocked_cyclic_sector_derived_properties}
    Start from
    Theorem~\ref{thm:ft_after_blocking_per_block_cyclic_nonzero_blocks}. With
    $p_A,p_B$ the least common multiples of the period-removal lengths on the two
    sides, use $p=\lcm(p_A,p_B)$. The prescribed-common-multiple
    lemma constructs both one-sided sector families at this length, and the
    positive-length blocking lemmas carry the one-sided tensor/nonzero-part
    equalities over to the common alphabet; for every $N\ge 1$,
    \[
        V^{(N)}\!(A^{[p]})_\sigma
        =
        V^{(N)}\!\bigl(\textstyle\bigoplus_j
        (\mu^A_j)^p (A_j)^{[p]}\bigr)_\sigma.
    \]
    The same applies on the $B$ side, and the two-sided blocking lemma gives
    equality of the two powered nonzero parts.
\end{proof}

\section{Blocking and weighted direct sums}%
\label{subsec:internal_lean_only_labels}

Blocking by a length~$p$ commutes with weighted direct sums up to the $p$th power
of each weight: for nonzero weights $(\mu_k)$ and blocks $(B_k)$,
\[
    \left(\textstyle\bigoplus_k \mu_k B_k\right)^{[p]}
    \sim \textstyle\bigoplus_k \mu_k^{\,p}\, B_k^{[p]},
\]
where $\sim$ is equality of MPV families. When two block families are reblocked to
a common length $p = m_k e_k$, the $p$-blocked alphabet is identified with the
$e_k$-fold tensor power of the $m_k$-blocked alphabet, so each sector reblocked
through its own period agrees with the directly $p$-blocked block.

\begin{lemma}[Adjoint primitivity equivalence]%
    \label{thm:primitive_adjoint_equiv}
    \lean{MPSTensor.IsPrimitive.adjoint_iff}\leanok
    \uses{def:primitive_channel}
    Let $\E$ be a finite-dimensional completely positive map. Then $\E$ is
    primitive if and only if its Frobenius adjoint $\E^\dagger$ is primitive;
    equivalently, $\sigma_\partial(\E) = \{1\}$ iff
    $\sigma_\partial(\E^\dagger) = \{1\}$, since the peripheral spectrum of
    $\E^\dagger$ is the complex conjugate of that of $\E$.
\end{lemma}

\begin{lemma}[Primitivity under blocking by a multiple]%
    \label{thm:primitive_blocking_divisible}
    \lean{MPSTensor.isPrimitive_transferMap_blockTensor_of_dvd}\leanok
    \uses{def:blocked_tensor, def:primitive_channel}
    Let $A$ have positive bond dimension and let $p_0,p\ge 1$ with $p_0 \mid p$.
    If $\E_{A^{[p_0]}}$ is primitive, then $\E_{A^{[p]}}$ is primitive.
\end{lemma}

\begin{lemma}[Blocking the weighted block sum]%
    \label{thm:block_tensor_to_tensor_from_blocks}
    \lean{MPSTensor.sameMPV₂_blockTensor_toTensorFromBlocks}\leanok
    \uses{def:blocked_tensor, def:block_diagonal_tensor, def:same_mpv2}
    For nonzero weights $(\mu_k)_{k=1}^r$, blocks $(B_k)_{k=1}^r$, and any
    $p\ge 1$,
    \[
        (\textstyle\bigoplus_k \mu_k B_k)^{[p]}
        \sim  \textstyle\bigoplus_k \mu_k^{\,p}\, B_k^{[p]},
    \]
    where $\sim$ denotes generation of the same MPV family at every length.
\end{lemma}

\begin{lemma}[Positive-length equality survives blocking]%
    \label{thm:same_mpv2_pos_block_tensor}
    \lean{MPSTensor.sameMPV₂Pos_blockTensor}\leanok
    \uses{def:blocked_tensor, def:same_mpv2_pos}
    If two tensors have the same MPV coefficients at every positive length, then
    so do their blocked tensors after any positive blocking length.
\end{lemma}

\begin{lemma}[Positive-length blocking of a weighted block sum]%
    \label{thm:same_mpv2_pos_block_tensor_to_tensor_from_blocks}
    \lean{MPSTensor.sameMPV₂Pos_blockTensor_toTensorFromBlocks}\leanok
    \uses{def:blocked_tensor, def:block_diagonal_tensor, def:same_mpv2_pos,
        thm:same_mpv2_pos_block_tensor, thm:block_tensor_to_tensor_from_blocks}
    If $A$ and $\bigoplus_k \mu_k B_k$ have the same MPV coefficients at every
    positive length, then after any positive blocking length $p$, $A^{[p]}$ and
    $\bigoplus_k \mu_k^{\,p} B_k^{[p]}$ have the same MPV coefficients at every
    positive length.
\end{lemma}

\begin{lemma}[Positive-length equality of powered block sums]%
    \label{thm:same_mpv2_pos_to_tensor_from_blocks_block_power}
    \lean{MPSTensor.sameMPV₂Pos_toTensorFromBlocks_blockPower}\leanok
    \uses{def:blocked_tensor, def:block_diagonal_tensor, def:same_mpv2_pos,
        thm:same_mpv2_pos_block_tensor, thm:block_tensor_to_tensor_from_blocks}
    If $\bigoplus_j \mu^A_j A_j$ and $\bigoplus_k \mu^B_k B_k$ have the same MPV
    coefficients at every positive length, then after any positive common
    blocking length $p$, $\bigoplus_j (\mu^A_j)^p A_j^{[p]}$ and
    $\bigoplus_k (\mu^B_k)^p B_k^{[p]}$ have the same MPV coefficients at every
    positive length.
\end{lemma}

\begin{lemma}[Nonzero weights survive blocking]%
    \label{lem:block_weights_ne_zero}
    \lean{MPSTensor.blockWeights_ne_zero}\leanok
    If $\mu \in \C$ is nonzero and $p \ge 1$, then $\mu^p \neq 0$; blocking
    preserves the nonzero-weight hypothesis of a weighted block decomposition.
\end{lemma}

\begin{lemma}[Common reblocking of cyclic-sector families]%
    \label{thm:exists_common_blocked_cyclic_sector_family}
    \lean{MPSTensor.exists_commonBlockedCyclicSectorFamily_of_hasPrimitiveIrreducibleCyclicSectors}\leanok
    \uses{def:primitive_irreducible_cyclic_sectors, def:blocked_tensor}
    Given a finite family $(B_k)_{k=1}^r$, each with primitive irreducible cyclic
    sectors, there is a common blocking length $p$ such that each $B_k^{[p]}$
    admits a unit-weight direct-sum decomposition over a single common blocked
    alphabet $d^{[p]}$ into left-canonical blocks with primitive transfer maps,
    tensor irreducibility, and positive bond dimensions.
\end{lemma}

\begin{lemma}[Sector irreducibility for the adjoint blocked map]%
    \label{thm:sector_irred_unconditional}
    \lean{MPSTensor.isIrreducibleOnCorner_of_cyclic_decomp_mps}\leanok
    \uses{def:irreducible_tensor}
    Let $A$ be a left-canonical irreducible MPS tensor with $D>0$, let $m$ be the
    period of its cyclic-sector decomposition, and let $(P_u)_{u=1}^m$ be the
    cyclic projections of the adjoint transfer map with
    $\E_{A^\dagger}(P_{u+1}) = P_u$ and $\sum_u P_u = \Id$. Then the restriction
    of $(\E_{A^\dagger})^m$ to each corner $P_u \MN{D} P_u$ is irreducible.
\end{lemma}

\begin{lemma}[Iterated blocking at the transfer-map level]%
    \label{thm:iterated_blocking_transfer}
    \lean{MPSTensor.transferMap_blockTensor_mul}\leanok
    \uses{def:blocked_tensor, def:transfer_map}
    For any MPS tensor $A$ and $m,n \ge 0$,
    \[
        \E_{(A^{[m]})^{[n]}} =  \E_{A^{[mn]}}.
    \]
\end{lemma}

\begin{definition}[Primitive irreducible cyclic-sector decomposition]%
    \label{def:primitive_irreducible_cyclic_sectors}
    \lean{MPSTensor.HasPrimitiveIrreducibleCyclicSectors}\leanok
    \uses{def:blocked_tensor, def:irreducible_tensor, def:primitive_channel,
        def:tp_kraus}
    $A$ \emph{has primitive irreducible cyclic sectors} of period $m \ge 1$ if
    there is a family $(C_k)_{k=1}^m$ with $A^{[m]} \sim \bigoplus_{k=1}^{m} C_k$
    (unit weights), each $C_k$ left-canonical, with primitive transfer map,
    tensor-irreducible, and of positive bond dimension.
\end{definition}

\begin{lemma}[Common reblocking at a prescribed multiple]%
    \label{thm:exists_common_blocked_cyclic_sector_family_common_multiple}
    \lean{MPSTensor.exists_commonBlockedCyclicSectorFamily_of_commonMultiple}\leanok
    \uses{thm:exists_common_blocked_cyclic_sector_family}
    With hypotheses as in
    Lemma~\ref{thm:exists_common_blocked_cyclic_sector_family}, and given any
    common multiple $p$ of the per-block period-removal lengths $(m_k)$, the
    common reblocking can be performed at the prescribed length $p$.
\end{lemma}

\begin{lemma}[Structural properties of the common reblocked sectors]%
    \label{thm:common_blocked_cyclic_sector_derived_properties}
    \lean{MPSTensor.CommonBlockedCyclicSectorFamily.derived_properties}\leanok
    \uses{thm:exists_common_blocked_cyclic_sector_family}
    Each sector tensor in the common reblocked cyclic-sector family of
    Lemma~\ref{thm:exists_common_blocked_cyclic_sector_family} is
    trace-preserving ($\sum_i (B^i)^\dagger B^i = \Id$), has primitive transfer
    map ($\sigma_\partial(\E_B) = \{1\}$), is tensor-irreducible, and has positive
    bond dimension.
\end{lemma}
